\section{Data and corpus construction}
\label{sec:data}

The corpus comprises 343 unique corporate reports from 49 companies
headquartered in seven European countries (Belgium, Finland, France,
Germany, Italy, the Netherlands, Spain), one report per company per year
from 2018 to 2024. The panel is exactly balanced: every year contributes 49
documents, and at a finer grain every country contributes a fixed number of
documents in every year of the panel (France 17, Germany 14, the
Netherlands 6, Italy 5, Spain 4, Finland 2, Belgium 1). Balance here is
temporal, not cross-national: countries are represented in wildly unequal
numbers, as Table~\ref{tab:data-country} shows, but each country's own
count is identical in every one of the seven years, so no year departs
compositionally from any other.

Three document types are represented: annual reports, the Universal
Registration Document (URD; French \emph{document d'enregistrement
universel}), and standalone sustainability reports. Document type was assigned per report from an
external metadata file (\texttt{muestras\_informes.xlsx}), joined to each
PDF by country, publication year and a normalised company name --- Unicode-
folded, punctuation-stripped, and reconciled through a manual company-name
map covering divergences between the directory tree and the metadata
source. The join resolved all 343 documents; no case was left unmatched.

\subsection{The country/type confound is a structural property of the corpus}

Table~\ref{tab:data-country} reports the corpus by country and document
type. Document type is confounded with country almost completely: 99 of the
corpus's 101 URD filings are French, and all 98 German documents are annual
reports. Both are properties of national reporting practice rather than
sampling accidents --- the URD is the standard annual regulatory filing for
French listed companies and routinely carries the extended sustainability
disclosure, whereas the German issuers in this sample report ESG content
inside the ordinary annual report. We state this here as a design property
of the corpus, not as a limitation discovered later: within the pooled
panel, a country effect and a document-type effect are not separately
identifiable. Because Germany is simultaneously a single country and a
single document type (98 of 98 documents are German annual reports), it is
the one slice of the panel where a temporal estimate is free of this
confound, and \S\ref{sec:temporal} reports a within-Germany regression for
that reason.

\begin{table}[htbp]
\centering
\caption{Corpus composition by country and document type.}
\label{tab:data-country}
\begin{tabular}{lrrrrr}
\toprule
Country & Companies & Annual report & URD & Sustainability report & Total \\
\midrule
France           & 17 & 13  & 99  & 7 & 119 \\
Germany          & 14 & 98  & 0   & 0 & 98  \\
Netherlands      & 6  & 40  & 2   & 0 & 42  \\
Italy            & 5  & 35  & 0   & 0 & 35  \\
Spain            & 4  & 28  & 0   & 0 & 28  \\
Finland          & 2  & 14  & 0   & 0 & 14  \\
Belgium          & 1  & 7   & 0   & 0 & 7   \\
\midrule
Total            & 49 & 235 & 101 & 7 & 343 \\
\bottomrule
\end{tabular}
\end{table}

Table~\ref{tab:data-year} reports the same corpus by year. Composition by
document type is stable across the panel (annual reports 33--36 per year,
URD 12--15, sustainability reports 1 per year), which matters when the
temporal series (\S\ref{sec:temporal}) is read against a fixed compositional
backdrop rather than a shifting one.

\begin{table}[htbp]
\centering
\caption{Corpus composition by year.}
\label{tab:data-year}
\begin{tabular}{lrrrr}
\toprule
Year & Annual report & URD & Sustainability report & Total \\
\midrule
2018 & 36  & 12  & 1 & 49  \\
2019 & 33  & 15  & 1 & 49  \\
2020 & 33  & 15  & 1 & 49  \\
2021 & 33  & 15  & 1 & 49  \\
2022 & 34  & 14  & 1 & 49  \\
2023 & 33  & 15  & 1 & 49  \\
2024 & 33  & 15  & 1 & 49  \\
\midrule
Total & 235 & 101 & 7 & 343 \\
\bottomrule
\end{tabular}
\end{table}

\subsection{Comparability across document types}

The panel spans three document types with different structure and
different length: annual reports in which ESG disclosure is one section
among many financial and operational ones, URD filings that combine
regulatory financial, governance and non-financial statements, and
standalone sustainability reports that are ESG content throughout. Raw
comparability across types cannot be assumed. Every measurement from
\S\ref{sec:extraction} onward accordingly operates not on whole documents but
on the ESG-relevant zones the lexical filter identifies within each one,
which is what lets a single index construction (\S\ref{sec:index}) be applied
uniformly across all three types. \S\ref{sec:scope} sets out how this corpus
relates to the source study's; here it is the object of measurement.

\subsection{Deduplication}

Corpus assembly began from 344 collected PDF filings, one more than the
final corpus size. A content-hash comparison identified a single duplicate
pair: a 2019 annual report for Wolters Kluwer (Netherlands) and a
\texttt{\_repaired} re-issue of the same filing under a different filename.
The hash is taken over the preprocessed representation of
\S\ref{sec:lexicon} rather than the raw extraction, so the criterion is
that two documents are indistinguishable to every downstream measure; here
the two files are in fact identical at the extraction stage as well. The
\texttt{\_repaired} copy was removed as a data-quality control, leaving 343
unique reports. The balanced panel reported above --- 49 documents in every
year, and a fixed per-country count within every year --- holds only after
this step: before deduplication, 2019 carried 50 documents, with the
Netherlands contributing 7 rather than 6, while every other year already
stood at 49.
